\section{Recurrent amplification and numerical checks}\label{supp:amplification}
These protocols support the coupling, timing, and field-induced gamma results in
\S\ref{sec:jcurve}--\S\ref{sec:qif}.
\subsection{JR coupling-sweep settings}\label{app:jr-sweep-settings}
The coupling sweeps of Figs.~\ref{fig:jcurve} and~\ref{fig:timing}, and the transcranial
alternating-current stimulation (tACS) control of Fig.~\ref{fig:tacs}, use the single JR
column with the field-free operating point $v^*$ pinned: for
each global connectivity $C$ we solve the field-free fixed-point relation analytically
for the external drive $p(C)$ that holds $v^*$ (hence $\Sig''(v^*)$) fixed, vary $C$
toward the alpha Hopf, and read $\gamma,\omega_0$ from the $6\times6$ Jacobian. The field
is the AM surrogate of \ifsupp main-text\fi\ Eq.~\eqref{eq:amsur} for Fig.~\ref{fig:jcurve}, genuine two
tones for the timing comparison (Fig.~\ref{fig:timing}), or a direct in-band sinusoid $s(t)=\varepsilon\cos\Omega t$
($\varepsilon=0.1$~mV) for the tACS control (Fig.~\ref{fig:tacs}), each read out by the
lock-in of $v$ at $\Delta f$, except in the timing comparison, which reads the
modulation and mean of the same pyramidal rate $r_P=\Sig(v+s)$.
That comparison uses $18$~s settling and a $20$~s window for each waveform.
The Figure~\ref{fig:jcurve}--Figure~\ref{fig:tacs} response-ratio comparison uses
$v^*=8.30$~mV, TI surrogate amplitude $0.3$~mV with $m=1$ and $f_c=100$~Hz,
and tACS amplitude $0.1$~mV, driven at each field-free modal frequency.
The original TI amplitudes are interpolated logarithmically in $C$ onto the tACS
grid. For the waveform comparison, both the surrogate and two-tone TI are rerun
on that grid with $12$~s settling and a $6$~s window.
Pinning fixes the background curvature; it does not cancel the stimulation-induced
mean shift.
Additional controls at $C=110$, $133.9091$, and $136.3$ use a tACS amplitude of
$0.0143039$~mV, matched to the leading TI beat source. They record mean voltage,
mean rate, and lock-in with $12$~s settling and a $6$~s window, repeated with
$60$~s settling and a $20$~s window; the latter settings also test genuine two tones.


\ifjne
\subsection{PING coupling-sweep settings}\label{supp:ping-runs}
We compute the linear reference's
second-order source from the PING quadratic terms at the central carrier,
$300$~Hz, and apply the full linearized network response over $15$--$90$~Hz
in $0.025$~Hz steps. Using the distinct AM sidebands changes the peak gain by
less than $0.1\%$ across this sweep.

The simulated coupling sweep uses the AM surrogate, $A_f=8$, with $5$~s
settling and a $4$~s window. We search around each linear resonance, extending
the frequency grid until each sampled maximum is bracketed, and compare
amplitudes $A_f=2,4,8$ at nine couplings. Local parabolic interpolation changes
the sampled peaks by less than $1.3\%$. The figure scales the linear reference
to the median measured gain over $14.5<J<16$; the amplitude check instead
normalizes peak response divided by $A_f^2$ and linear gain to its value at
$A_f=2$, $J=15.53$. These normalizations compare gain variation with the detector coefficient
allowed to vary along this unpinned sweep.


\fi

\subsection{Stability of the driven PING endpoint}\label{app:ping-stability}
We check the stability of the driven $J$-curve endpoint by integrating the full
variational matrix along its trajectory. At carrier $300$~Hz and beat $50.5$~Hz,
the AM input repeats every $2$~s. After $8$~s of settling, we compute the
perturbation multipliers over that period and repeat over the next period.
Moduli below one indicate decay of perturbations to the periodic response.
Two-tone controls at $J=16.5663$ and $16.5901$, both with a $50.5$~Hz beat,
use $5$~s settling, direct perturbation decay, and an $8$~s finite-time
Lyapunov estimate.

Although a DC-only averaged calculation predicts an unstable endpoint, the full
periodically forced solution remains attracting: its largest perturbation
multiplier is $0.034$ over the $2$~s forcing period. Increasing settling from
$0.6$ to $3$~s changes the endpoint lock-in by $0.8\%$. With genuine two tones,
the $50.5$~Hz endpoint control also attracts nearby trajectories: perturbations
decay at $-0.0014$~ms$^{-1}$, and the $8$~s Lyapunov estimate is
$-0.00079$~ms$^{-1}$. These controls test a nonlinear driven response, not linear
gain about the DC-shifted equilibrium.

\subsection{Finite-network numerical checks}\label{app:qif-checks}
At $A_f=10$ and $\bar\eta=0.85$, the finite-network beat depth is $0.093$
for the published cutoff, versus the mean-field value $0.128$. Raising the spike cutoff from $V_{\rm peak}=100$ to
$200$ and $400$ gives depths $0.127$ and $0.139$, while the mean-rate increase
stays at $2.8$--$2.9\%$. At the original cutoff, doubling the population gives
$0.083$ and halving the time step gives $0.093$. The response is therefore
sensitive to the finite spike cutoff. Halving the step at cutoffs $200$ and
$400$ gives depths $0.121$ and $0.141$, respectively: cutoff dependence persists,
and quantitative convergence remains unestablished at this population size.
At the plotted cutoff, the pooled spike-phase vector strength is $0.046$;
median per-neuron phase consistency remains negative in field-off and field-on
windows at all three cutoffs.
In the $A_f=35$ convergence checks, halving the time step changes the lock-in and the pooled vector
strength by under $1\%$; doubling the spike-detection threshold changes the lock-in by
about $1\%$ and the mean rate by $2\%$, but the pooled vector strength at
$\bar\eta=0.85$ by about $45\%$, so phase concentration, unlike rate modulation, is sensitive to
how spike times are defined.


\ifjne
\subsection{A cutoff-free theta-neuron check}\label{supp:theta-pilot}
The change of variables $V_j=\tan(\theta_j/2)$ converts the QIF equation to
\[
\tau_X\dot\theta_j=(1-\cos\theta_j)
 +(1+\cos\theta_j)(\eta_j+U_X(t)),
\]
where $U_X$ is the same recurrent and external input as in the voltage model.
An upward crossing of $\theta=\pi$ emits a spike; wrapping by $2\pi$ replaces
the finite voltage cutoff and reset. We linearly interpolate the crossing
time for phase statistics. Empirical spike counts still drive the original
second-order synapses. Excitability quantiles, initial voltage, coupling,
carrier, and readouts are unchanged.

The pilot used $\bar\eta=0.85$, $A_f=10$, $f_c=300$~Hz, a $55$~Hz beat,
and $800$~ms settling. Before running it, we set a stopping rule: rate depth
must change by less than $5\%$ under separate step halving, population
doubling, and window doubling before extending this solver to carrier controls.
Table~\ref{tab:theta-pilot} gives every tested setting. Step and window changes
satisfy that criterion; population doubling does not, so the pilot stops.
No finite-network equal-RMS control was run.

\begin{table}[htbp]
\centering\small
\begin{tabular}{@{}lrrrrr@{}}
\toprule
Condition & $N$ per population & Step (ms) & Window (ms) & Mean E rate (Hz) & Beat depth\\
\midrule
Field off & 8000 & 0.0025 & 800 & 50.4 & 0.006\\
Two tones & 8000 & 0.0025 & 800 & 51.8 & 0.156\\
Half step & 8000 & 0.00125 & 800 & 51.8 & 0.155\\
Double population & 16000 & 0.0025 & 800 & 51.8 & 0.137\\
Double window & 8000 & 0.0025 & 1600 & 51.8 & 0.151\\
Exact mean field & $\infty$ & 0.0025 & 800 & 52.1 & 0.128\\
\bottomrule
\end{tabular}
\caption{Cutoff-free weak-drive checks. Depth is the rate lock-in divided by
mean rate. Relative to the two-tone baseline it changes by $-0.4\%$,
$-11.9\%$, and $-3.0\%$ under step, population, and window changes.
Only the baseline settings have a matched field-off run.}
\label{tab:theta-pilot}
\end{table}

Removing the cutoff raises the baseline depth from $0.093$ to $0.156$.
The three cutoff results are compatible with an inverse-cutoff correction:
a fit of $d_\infty-c/V_{\rm peak}$ gives $d_\infty=0.156$, with residuals
below $0.002$. This agreement supports a cutoff contribution but does not
establish joint numerical convergence. Doubling the population moves the
theta result toward the exact mean field, while also increasing the largest
sampled Euler step in phase from about $1.5$ to $3.2$ radians. A combined
population and time-step limit remains untested.

Pooled spike-phase vector strength rises from $0.002$ without stimulation
to $0.068$--$0.077$ in the driven runs. Median per-neuron phase consistency
remains negative, from $-0.015$ to $-0.005$, and supplies no evidence for
strong individual phase locking. The finite network therefore supports the
presence of population-rate modulation; its quantitative weak-input depth
and agreement with the mean field remain unresolved.

\fi
\subsection{Fast-carrier estimate at the illustrated amplitudes}
The $A_f=35$ runs are finite-amplitude illustrations: the fast-carrier estimate gives a voltage excursion
$A_f/(\tau_E\omega_c)\approx1.2$ per tone in the model's voltage units. Its
rectified excitatory-drive shift is about $1.5$, compared with $0.12$ at
$A_f=10$. At $\bar\eta=0.85$, the Hopf threshold for a drive shift applied
only to the excitatory population is $\delta\eta_E^*\approx0.286$. The larger
field is therefore expected to cross this threshold, whereas the smaller one
remains below it.

\setcounter{figure}{2}
\begin{figure}[t!]
\centering
\includegraphics[width=0.92\textwidth]{fig_tacs_jcurve.pdf}
\caption{\textbf{tACS shares the network amplifier (Jansen--Rit).} A direct
low-frequency (tACS) drive at $\Delta f$ enters the column without a carrier, so it
needs no demodulation, but feeds the same near-Hopf resonator. With the field-free operating point
pinned, sweeping the coupling $C$ toward the alpha Hopf: \emph{(a)} the lock-in
resonance at $\Delta f$ sharpens and grows as $C\!\to\!C^*$; \emph{(b)} the response
grows as $1/\gamma$ (dashed) and saturates near criticality---the same network
amplification as the demodulated TI drive (Fig.~\ref{fig:jcurve}). The leading tACS
response is linear in the field; the TI beat is quadratic. Their response ratio
is independent of damping in linear response and changes at the finite amplitudes
nearest the Hopf (\S\ref{sec:jcurve}). Here $\gamma>0$ describes the field-free state.}
\label{fig:tacs}
\end{figure}

\setcounter{figure}{3}
\begin{figure}[t!]
\centering
\includegraphics[width=\textwidth]{fig_entrainment.pdf}
\caption{\textbf{Frequency dominance in the driven autonomous Jansen--Rit column.}
The column is set in its autonomous alpha limit cycle ($p=250$, $f_0\!\approx\!11$~Hz) and
driven by a direct in-band sinusoid. \emph{(a)} The set
$|f_{\rm out}-\Delta f|<0.3$~Hz forms a wedge near $f_0$ that widens with drive.
It screens for phase capture but does not establish it. \emph{(b)} Inside this
region the dominant frequency follows $\Delta f$; outside, it remains near $f_0$.
\emph{(c)} Toward the Hopf at fixed drive the dominance width grows as the
autonomous cycle amplitude vanishes (blue, left axis; red dashed, right axis). Above the
Hopf this width describes frequency selection in a driven autonomous oscillator;
it is distinct from the linear susceptibility of the stable-focus side
(Fig.~\ref{fig:jcurve}). Both diagnostics increase toward criticality in these sweeps.}
\label{fig:entrain}
\end{figure}

\setcounter{figure}{7}
\begin{figure}[t!]
\centering
\includegraphics[width=0.72\textwidth]{fig_nmm2_map.pdf}
\caption{\textbf{NMM2 PING resonance map (exact mean field).} Demodulated response over
the ($\Delta f$, background drive $\bar\eta$) plane, log color scale. The ridge follows the
gamma natural frequency, with $f_0$ rising with $\bar\eta$ for these parameters. White markers
outline the frequency-dominance set (\S\ref{sec:vocab}), drawn only above the Hopf
(dashed line) where there is an autonomous rhythm to capture; below it the same color field
is a forced response. The markers do not establish phase capture. Lock-in amplitude
cannot separate driven and autonomous contributions, since both can place power at
$\Delta f$.}
\label{fig:nmm2_map}
\end{figure}

\setcounter{figure}{9}
\begin{figure}[t!]
\centering
\includegraphics[width=0.92\textwidth]{fig_nmm2_jcurve.pdf}
\caption{\textbf{Coupling changes the dominant mode and its gain in NMM2.}
Unpinned sweep at $\bar\eta=0$, AM surrogate $A_f=8$.
\emph{(a)} Simulated rate peak (blue) and linear reference including the
second-order source (gray), normalized over $14.5<J<16$. Red dashed:
field-free Hopf. \emph{(b)} Gain against inverse damping of the critical
mode. Filled symbols: that mode carries the linear peak; open symbols:
the competing $55$~Hz mode carries it ($14.6\lesssim J\lesssim16.2$).
The dashed $1/\gamma$ guide is anchored at the last linear point.
The simulation exceeds then falls below the guide. The driven endpoint
remains attracting under both AM and genuine two-tone forcing.}
\label{fig:nmm2_jcurve}
\end{figure}

\setcounter{figure}{10}
\begin{figure}[t!]
\centering
\includegraphics[width=\textwidth]{fig_timing_not_rate.pdf}
\caption{\textbf{Beat modulation grows more than the mean-rate shift.}
JR coupling sweep with fixed field-free curvature and genuine two-tone
input, $\Delta f=f_0$. Both observables use the pyramidal rate
$r_P=\Sig(v+s)$, in mHz.
\emph{(a)} Beat amplitude (blue) follows $1/\gamma$ (dashed) before
saturating; mean-shift magnitude (red) remains smaller.
\emph{(b)} Normalization to the far-from-Hopf values gives growth factors
$13.2$ and $4.4$. Their ratio rises from $14.5$ to $58$ near
$\gamma=1$~s$^{-1}$, then falls to $44$ as the amplified response rectifies.
The AM surrogate gives ratios $19$, $66$, and $45$.}
\label{fig:timing}
\end{figure}

\setcounter{figure}{11}
\begin{figure}[t!]
\centering
\includegraphics[width=\textwidth]{fig_qif_timing_twotone.pdf}
\caption{\textbf{Beat-locked modulation and mean-rate change in the QIF network.} \emph{(a)} Field-free power
spectrum of $r_E(t)$, normalized: the spiking network ($N=8000$) matches the exact NMM2
mean field in gamma frequency ($\sim$55~Hz); a normalized spectrum does not test the mean
rate. \emph{(b)} Intermodulation above the
Hopf (TI off vs.\ on, log power): the nonlinearity mixes the imposed $\Delta f=42$~Hz beat
with the persisting gamma near $53$--$54$~Hz, raising the direct $\Delta f$ line and the sum sideband
near $95$~Hz, both absent with TI off. These sidebands establish nonlinear mixing;
they do not uniquely identify gamma-amplitude gating. \emph{(c)} The $\Delta f$ modulation depth
$A_\Omega/\langle r_E\rangle$, off vs.\ on: below $0.003$ rising to $0.29$
($\bar\eta=0.85$) and $0.23$ ($\bar\eta=1.5$), while the mean rate rises by about
$8$--$11\%$. The labels identify field-free states; both field-on responses contain
gamma activity beyond the beat component. Spectra use the same post-settling windows
as the rate measurements and are normalized to each spectrum's maximum. Modulation depth is
reported rather than the off$\to$on lock-in ratio, whose denominator is the field-free
floor.}
\label{fig:qif_timing}
\end{figure}

